% L3-2 revised (2026-10-28, far-field refinement L3-2r): replaces \subsubsection{Point resonator on the plate}
% in §6.5 of numerical_results.tex (label subsubsec:num_L3_2), figure environment included.
% Figure: plots/L3-2/point_resonator.png -> figures/L3_2_point_resonator.png (scripts/export_paper_assets.jl, unchanged mapping).
% Data: data/exp_pro/L3-2/{points,convergence,fits,winkler,checks}.csv (frozen ladder L1-L3: 136/136; added L2-L4 rows: 42/56);
% first run (graded far field, 113/136): data/exp_pro/L3-2/graded_far_field/.
{\color{red}
\subsubsection{Point resonator on the plate}\label{subsubsec:num_L3_2}

\paragraph{Test description.} A point resonator is attached to the plate in a channel of width $W=2~\si{\meter}$ and depth $H=10~\si{\meter}$, with walls that carry no flux. Along the channel there is a damping zone of $4\lambda_0$, open water of $\lambda_0/2$, the plate, open water and a second zone. The plate is $4\lambda_p$ long, spans the width with free lateral edges, and has $C_\rho=0.05g$, $m_\rho=0.01~\si{\meter}$ and $\nu=0$. The resonator sits at the plate centre $\xvec_r$, with $M_\rho=0.01~\si{\meter\cubed}$ and $\omega_r=10~\si{\radian\per\second}$. An incident wave enters from the inlet. The fields are $Q_2$. On the plate the far-field cells are fixed at $\lambda_p/16$ along and $W/14$ across the channel, the same on every mesh. Only the two far-field cells next to $\xvec_r$, in each direction, are replaced by cells graded from $h_\mathrm{min}=\lambda_p/16$, $\lambda_p/32$, $\lambda_p/64$ or $\lambda_p/128$ (ratio at most $1.5$; $N=1.3\times10^5$ to $2.1\times10^5$). The Kirchhoff Green function is finite at the source, with $\eta-\eta(\xvec_r)\sim c\,r^2\ln r$ and $c=f/(8\pi C_\rho)$, so no footprint is needed. We solve the bare plate under the incident wave and under a unit point load $f$, which gives $G(\xvec_r,\xvec_r)$, and the plate with the resonator. The attachment point is a mesh node, the midpoint of a cell edge or a cell interior point. Two softer plates, $C_\rho=0.005g$ and $5\times10^{-4}g$, are run in a channel scaled with $\lambda_p$. The static Winkler plate, $C_\rho\nabla^4w+gw=f\delta$, gives an independent check of the singularity. The sparse LU factorisation uses a nested-dissection (METIS) ordering. Every run stays below $10$~GB of memory.

\paragraph{Results.} The resonance identities hold to round-off at every attachment point and every $C_\rho$. At $\omega_r$ with $\zeta=0$, $\eta_h(\xvec_r)$ vanishes to $1.4\times10^{-15}$ of the bare value. The Sherman--Morrison relation between $q_h$, the bare field and $G_h$ holds to $1.2\times10^{-9}$ at $\omega_r$ and to $3.6\times10^{-9}$ at $0.9\omega_r$ with $\zeta=0.02$. With the far field fixed, the driving point converges (Figure~\ref{fig:num_L3_2}b). Over the first three meshes, $G_h(\xvec_r,\xvec_r)$ converges at rates of $1.6$ to $1.9$ and $|q_h|$ at $1.8$ to $2.1$, for all three plates. The changes between the second and third meshes are $0.15$--$0.24\%$ for $G_h$ and $0.03$--$0.12\%$ for $|q_h|$ (tolerance $2\%$). The near field around $\xvec_r$ converges at rates of $1.5$ in $L^2$ and $1.7$ in $H^1$. On the finest mesh the changes level off: $G_h$ moves by about $0.3\%$ for every plate, and this does not decrease. The same increment appears with a top cell half as thick or with 12 vertical cells, so it does not come from the vertical mesh. Most likely it comes from the elongated cells that the tensor-product grading creates along the two lines through $\xvec_r$ (aspect ratio $8$ on the finest mesh). The driving-point compliance is therefore resolved to about $0.3\%$, against about $5\%$ in a first ladder whose far-field cells changed with the refinement. The local field has the predicted singularity (Figure~\ref{fig:num_L3_2}a). On the third mesh the fitted coefficient gives $c_h\,8\pi C_\rho/f_h=1$ within $2.3\%$ for the resonator force and within $3.7\%$ for the point load, and within $3.9\%$ for the softer plates (tolerance $10\%$). The static plate gives $w_h(0)$ within $0.67\%$ of $f/(8\sqrt{gC_\rho})$, with a rate of $1.03$. The attachment location does not matter: $G_h$ at the edge and interior points differs from $G_h$ at the node by $0.06\%$ and $0.11\%$ on the third mesh, and the difference falls fourfold per refinement. The penalty control changes $G_h$ by $0.15\%$. The driving-point compliance stays bounded as $C_\rho$ decreases. It grows as $C_\rho^{-0.49}$, against $C_\rho^{-0.62}$ for the infinite plate on the fluid (Figure~\ref{fig:num_L3_2}c). We conclude that a point resonator on the plate needs no footprint, and that its driving-point response converges under local refinement.

\begin{figure}[t]
    \centering
    \includegraphics[width=\linewidth]{figures/L3_2_point_resonator.png}
    \caption{Point resonator on the plate in a three-dimensional channel, with the far field fixed at $\lambda_p/16$. (a) Azimuthal mean of $\mathrm{Re}\,\eta$ about the attachment point minus its centre value, third mesh ($h_\mathrm{min}=\lambda_p/64$), $C_\rho=0.05g$; markers: finite element; line: $c\,r^2\ln r$ with $c=f_h/(8\pi C_\rho)$ plus the fitted regular terms. (b) Differences of $G_h(\xvec_r,\xvec_r)$ and $|q_h|$ between successive meshes, relative to the finest value, against the $h_\mathrm{min}/\lambda_p$ of the finer mesh of each pair. (c) $|G_h(\xvec_r,\xvec_r)|$ on the third mesh against $C_\rho$; crosses: infinite plate on the fluid; dashed: $C_\rho^{-1/2}$.}
    \label{fig:num_L3_2}
\end{figure}
}
